\documentclass[12pt,a4paper]{article}

\usepackage[utf8]{inputenc}
\usepackage[brazil]{babel}
\usepackage{graphicx}
\usepackage{geometry}
\usepackage{tcolorbox}
\usepackage{tikz}
\usepackage{fancyhdr}
\usepackage{setspace}

\geometry{margin=2.9cm}
\setstretch{1.25}

\setlength{\headheight}{40pt}

\pagestyle{fancy}

\fancyhf{}

\fancyhead[L]{Manual Didático}
\fancyhead[C]{\includegraphics[height=1.7cm]{logo-ifmt.png}}
\fancyhead[R]{Projeto Semáforo Arduino}

\renewcommand{\headrulewidth}{0.4pt}

\fancyfoot[L]{Instituto Federal de Educação, Ciência e Tecnologia de Mato Grosso $\vert$ Nicoly Canteiro}
\fancyfoot[R]{\hspace{-0.10cm}\thepage}

\renewcommand{\headrulewidth}{0.4pt}
\renewcommand{\footrulewidth}{0.4pt}

\tcbset{
colback=gray!10,
colframe=black,
arc=3mm
}

\begin{document}

\begin{center}

{\Large \textbf{MANUAL DIDÁTICO}}

\vspace{0.3cm}

Montagem de Semáforo com Arduino no Tinkercad

\vspace{0.3cm}

Extensionista: Nicoly Canteiro $\vert$ Professor do Projeto de Extensão: Professor da disciplina

\end{center}

\section*{1 Introdução}

\noindent
Neste projeto será construído um pequeno semáforo eletrônico utilizando Arduino.

\noindent
O Arduino funciona como o cérebro do circuito e controla três LEDs que representam as cores de um semáforo.

\begin{itemize}

\item Vermelho --- Parar
\item Amarelo --- Atenção
\item Verde --- Seguir

\end{itemize}

\noindent
O circuito será montado no simulador Tinkercad.

\section*{2 Ilustração de um Semáforo}

\begin{center}

\begin{tikzpicture}

\draw[fill=black!80] (0,0) rectangle (2,6);

\fill[red] (1,5) circle (0.5);
\fill[yellow] (1,3) circle (0.5);
\fill[green] (1,1) circle (0.5);

\end{tikzpicture}

\end{center}

\section*{3 Atalhos Importantes no Tinkercad}

\begin{tcolorbox}

S --- iniciar ou pausar a simulação

R --- girar o componente selecionado

E --- abrir a tela de programação

CTRL + C / CTRL + V --- copiar e colar componentes

F --- centralizar a área de trabalho

\end{tcolorbox}

\section*{4 Materiais Necessários}

\begin{itemize}

\item 1 Placa Arduino
\item 1 Protoboard
\item 3 LEDs (vermelho, amarelo e verde)
\item 3 Resistores de 220 ohms
\item Fios de Conexão

\end{itemize}

\section*{5 Função dos Componentes}

\begin{itemize}

\item Arduino: controla o circuito e executa o programa.
\item Protoboard: permite montar circuitos sem solda.
\item LED: dispositivo que emite luz quando recebe corrente elétrica.
\item Resistor 220 ohms: limita a corrente para proteger o LED.
\item Fios: conduzem energia entre os componentes.

\end{itemize}


\newpage
\section*{6 Passo a Passo da Montagem}

\noindent\textbf{\large Passo 1}

Coloque o Arduino e a protoboard na área de trabalho do Tinkercad.

\noindent\textbf{\large Passo 2}

Posicione três LEDs na protoboard:

\begin{itemize}

\item LED vermelho na linha H coluna 10 e 11
\item LED amarelo duas linhas abaixo
\item LED verde duas linhas abaixo

\end{itemize}

\noindent\textbf{\large Passo 3}

Conecte um resistor de 220 ohms em cada LED.

Uma ponta do resistor deve ir ao LED e a outra ao GND.

\noindent\textbf{\large Passo 4}

Realize as conexões com fios:

\begin{itemize}

\item Pino 13 do Arduino até o LED vermelho
\item Pino 12 até o LED amarelo
\item Pino 11 até o LED verde
\item Um fio preto conecta os resistores ao GND

\end{itemize}

\section*{7 Fluxo da Energia}

Arduino $\rightarrow$ fio $\rightarrow$ LED $\rightarrow$ resistor $\rightarrow$ GND

Quando o Arduino envia energia para um pino, o LED correspondente acende.

\section*{8 Código Arduino}

\begin{tcolorbox}

\begin{verbatim}

// Definição dos pinos
int vermelho = 13;
int amarelo = 12;
int verde = 11;

void setup() {
  pinMode(vermelho, OUTPUT);
  pinMode(amarelo, OUTPUT);
  pinMode(verde, OUTPUT);
}

void loop() {

  // Verde ligado
  digitalWrite(verde, HIGH);
  digitalWrite(amarelo, LOW);
  digitalWrite(vermelho, LOW);
  delay(5000); // 5 segundos

  // Amarelo ligado
  digitalWrite(verde, LOW);
  digitalWrite(amarelo, HIGH);
  digitalWrite(vermelho, LOW);
  delay(2000); // 2 segundos

  // Vermelho ligado
  digitalWrite(verde, LOW);
  digitalWrite(amarelo, LOW);
  digitalWrite(vermelho, HIGH);
  delay(5000); // 5 segundos
}

\end{verbatim}

\end{tcolorbox}

\section*{ }
\section*{9 Resultado}

\noindent
Se todas as conexões estiverem corretas:

\begin{itemize}

\item O LED verde acende primeiro
\item Depois o LED amarelo
\item Logo após o LED vermelho

\end{itemize}

\noindent
Isso simula o funcionamento de um semáforo eletrônico.

\end{document}